\section{Introduction}
\label{sec:intro}

Precise measurements of ionization and electroluminescence signals in high-pressure xenon are
essential for the development of next-generation rare-event detectors and for techniques such as
ion tagging ("Ba-tagging") that promise powerful background rejection in neutrinoless double beta
decay ($0\nu\beta\beta$) searches. The RITA (Radium Ion Tagging) setup was designed to explore the
physics and engineering challenges of producing, tagging, and transporting single ions in a
high-pressure xenon environment.

\begin{figure}[t]
    \centering
    \includegraphics[width=0.95\textwidth]{experiment_scheme.png}
    \caption{Conceptual schematic of the RITA alpha-tagging concept (panels a–c).
        Panel (a) shows the alpha-coated silicon detector (top view) where alpha decays produce a prompt signal in the silicon and eject a recoiling daughter into the xenon.
        Panels (b) and (c) show the \textit{per-ion} tagging principle: the silicon pulse is analyzed by a single-channel analyzer (SCA) which accepts events inside an energy region of interest (ROI) corresponding to the \isotope{Th}{228}-alpha.
        When an accepted alpha is observed, the gate voltage is switched to allow the selected daughter ion (e.g., Ra$^{++}$) to pass toward the downstream sensor. If the alpha falls outside the ROI, the gate remains biased to reject the ion.}
    \label{fig:setup_concept}
\end{figure}

The RITA apparatus was built with a specific downstream application in mind: the development and
characterization of single-ion capture methods and molecular sensing approaches relevant to
barium-tagging. Because radium and barium are chemically analogous, Ra$^{++}$ serves as an ideal
surrogate to test chemical interactions with fluorescent molecules, immobilization strategies, and
single-molecule fluorescence readout systems. Recent Density Functional Theory (DFT) computations
confirm this structural and chemical equivalence. Specifically, computational analysis of the
interaction between Ra$^{++}$ and the fluorescent bicolor indicator FBI-G1 demonstrates that the
aza-crown ether binds the radium cation with comparable thermodynamic stability to barium. This
binding is driven by both N-lone pair interactions and a strong $\pi$-cation interaction with the
para-disubstituted phenyl group, resulting in the necessary $\pi$-decoupling of the aromatic
systems (a detailed summary of these DFT computations is provided in Appendix X).

The core operating principle of RITA, illustrated in Figure \ref{fig:setup_schematic}, relies on an
alpha-emitting coating (\isotope{Th}{228}) on a silicon detector that provides a coincident tag.
The emitted alpha is stopped in the silicon, while the recoiling daughter ion is ejected into the
xenon gas. By selecting only the alpha energy window corresponding to \isotope{Ra}{224} recoils and
switching a gate bias on the timescale of the recoil transit, the system acts as an electronic mass
and identity filter, producing a narrow, event-by-event stream of Ra$^{++}$ ions.

However, before these ions can be effectively transported and captured by a downstream molecular
sensor, it is critical to understand their fundamental behavior upon entering the active volume.
The charge-state dynamics and neutralization probabilities of ions drifting in high-pressure xenon
are highly complex and carry direct implications for the viability of any tagging scheme (as
recently discussed in \cite{peralta_conde_a_charge-state_2026}). Furthermore, understanding the
electron-ion recombination of heavy recoiling nuclei in xenon provides fundamental data necessary
for signal characterization and background mitigation in contemporary direct dark matter
experiments. In noble gas TPCs, the ability to discriminate potential dark matter signals from
backgrounds relies primarily on the ratio of extracted ionization (S2) to primary scintillation
(S1). This ratio is fundamentally governed by the electron-ion recombination probability along the
particle track. Heavy nuclear recoils produce highly dense ionization topologies compared to
electron recoils, leading to distinct, energy- and field-dependent recombination dynamics.
Accurately modeling this charge yield—particularly for radiogenic backgrounds such as radon
daughters recoiling into the active volume—requires robust experimental calibration of the
recombination fraction across a wide range of applied drift fields.

When a daughter ion recoils into the xenon gas, it produces a dense ionization track. A significant
fraction of these ionization electrons recombine with the surrounding ions before they can be
extracted and converted into secondary scintillation (S2). This recombination reduces the detected
electron yield and depends heavily on the local ionization density and the applied drift field.
Characterizing this recombination is therefore a prerequisite for interpreting observed S2 signals,
understanding charge survival, and optimizing ion-extraction techniques.

While Figure \ref{fig:setup_concept} depicts a conceptual downstream sensor surface to illustrate
the per-ion delivery concept, the present RITA implementation utilizes a Photomultiplier Tube (PMT)
as the primary optical readout to study the bulk gas dynamics. This manuscript presents a
quantitative characterization of electron recombination for heavy recoils using a highly
constrained drift chamber operating at 1 bar with a Thick Gas Electron Multiplier (THGEM). We focus
specifically on recoil species produced by the \isotope{Th}{228} decay chain—notably the
\isotope{Ra}{224} nuclei—and analyze how the detected electron yield and S2 area saturation depend
on the applied drift and electroluminescence fields.

% In addition to the primary recombination measurements, this manuscript documents several supporting
% experimental efforts that help interpret the observed yields: (A) a gate transmission and
% collection‑efficiency study, and (B) a neutralization/collection measurement where recoiling radium
% species are physically collected and assayed by gamma spectroscopy to estimate charge‑state
% survival in vacuum versus xenon. These auxiliary experiments provide independent constraints on how
% many ions are expected to reach the extraction region and how much neutralization may reduce the
% number of extractable ions.

The rest of the paper is organised as follows. Section~\ref{sec:experiment} describes the RITA
apparatus, detectors and data acquisition. Section~\ref{sec:methods} summarises the analysis
pipeline used to extract S1/S2 windows, alpha energy maps and recoil S2 areas, and to compute the
recombination fraction. Section~\ref{sec:results} presents the calibration, S2 histograms and the
main recombination plots for the selected runs. Section~\ref{sec:discussion} discusses the
implications for ion tagging and compares our results with simple columnar recombination
expectations. Finally, Section~\ref{sec:conclusions} summarises our findings and outlines the next
experimental steps.
